\begin{table*}[ht]
    \centering
    \caption{Sensitivity of the $K=1$ meta-learning advantage to non-landslide (negative) sample selection, under spatial cross-validation (mean $F_1$ over three seeds). Target negatives are re-generated by random sampling outside exclusion buffers of 500\,m, 1\,km and 2\,km, and by a slope-matched strategy whose negatives reproduce the positives' slope distribution (removing the steep-vs-flat confound). The meta-learned prior (Reptile) reuses the unchanged source meta-model. Features are re-extracted from rasters, so absolute $F_1$ differs slightly from Table~\ref{tab:spatial_results}; the comparison across conditions is internally consistent. Absolute $F_1$ shifts with the sampling choice (consistent with \citealt{abdelkader2025nonlandslide}), but the meta-learning \emph{lift} stays positive throughout (range $+3.0$ to $+8.1$\,pp) and \emph{increases} under the hardest slope-matched condition. Source: \texttt{results/stats/negatives\_sensitivity.csv}.}
    \label{tab:neg_sensitivity}
    \footnotesize
    \begin{tabular}{llccc}
        \toprule
        \textbf{Target} & \textbf{Negative-sampling strategy} & \textbf{Indep.\ $F_1$} & \textbf{Reptile $F_1$} & \textbf{Lift (pp)} \\
        \midrule
        Huasco ($N=278$) & Random, 500\,m buffer (baseline) & 0.734 & 0.775 & $+4.1$ \\
         & Random, 1\,km buffer & 0.713 & 0.750 & $+3.8$ \\
         & Random, 2\,km buffer & 0.733 & 0.763 & $+3.0$ \\
         & Slope-matched & 0.654 & 0.709 & $+5.5$ \\
        \midrule
        Copiap\'o ($N=19$) & Random, 500\,m buffer (baseline) & 0.700 & 0.730 & $+3.0$ \\
         & Random, 1\,km buffer & 0.593 & 0.647 & $+5.4$ \\
         & Random, 2\,km buffer & 0.585 & 0.666 & $+8.1$ \\
        \bottomrule
    \end{tabular}
\end{table*}
